\section{Extended Problem Examples}
\label{app:examples}

We present six complete problem--solution pairs from \ecagbench{}, selected to illustrate the range of problem types, difficulty levels, and mathematical content.

\subsection*{A.1: Counterexample --- Infinite Disjoint Union (Level 2)}

\begin{problembox}[title={\ecagid{0024} --- Basic Concepts}]
\textbf{Problem:} Let $k$ be a field. Is the scheme $\coprod_{i=1}^{\infty} \Spec(k)$ affine? Justify your answer by identifying which property of affine schemes it fails to satisfy.

\medskip
\textbf{Solution (abbreviated):}
The scheme $X = \coprod_{i=1}^{\infty} \Spec(k)$ is not affine.
The key property that fails is \emph{quasi-compactness}: affine schemes are quasi-compact, but $X$ has the discrete topology on an infinite set, so the cover $\{U_i = \Spec(k)\}$ has no finite subcover.

Alternatively, if $X$ were affine, we would have $X \cong \Spec(R)$ for $R = \Gamma(X, \mathcal{O}_X) \cong \prod_{i=1}^{\infty} k$.
But $\Spec(\prod_{i=1}^{\infty} k)$ is a connected scheme (it has no non-trivial idempotents besides those coming from finite products), contradicting the fact that $X$ is disconnected with infinitely many connected components.

\medskip
\textbf{Type:} Counterexample \quad \textbf{Difficulty:} 2 \quad \textbf{Verification:} CAS (check quasi-compactness of global sections ring)
\end{problembox}

\subsection*{A.2: Example --- Hopf Algebra Structures in Char $p$ (Level 1)}

\begin{problembox}[title={\ecagid{0003} --- Basic Concepts}]
\textbf{Problem:} Let $k$ be a field of characteristic $p > 0$. Write down explicitly the Hopf algebra structures (comultiplication, counit, antipode) on $k[x]/(x^p - 1)$ for $\mu_p$ and on $k[x]/(x^p)$ for $\alpha_p$. Using the substitution $u = x - 1$ to identify the underlying algebras as $k[u]/(u^p)$, show that $\mu_p \not\cong \alpha_p$ as group schemes.

\medskip
\textbf{Solution (abbreviated):}
Under $u = x-1$, both algebras become $k[u]/(u^p)$, but the comultiplications differ:
\begin{align*}
\Delta_\mu(u) &= u \otimes 1 + 1 \otimes u + u \otimes u \quad (\text{from } \mu_p) \\
\Delta_\alpha(u) &= u \otimes 1 + 1 \otimes u \quad (\text{from } \alpha_p)
\end{align*}
Any $k$-algebra automorphism of $k[u]/(u^p)$ has the form $u \mapsto \lambda u + (\text{higher terms})$ with $\lambda \neq 0$, but $\Delta_\mu$ has a quadratic term $u \otimes u$ that cannot be eliminated by any such automorphism.

\medskip
\textbf{Type:} Example \quad \textbf{Difficulty:} 1 \quad \textbf{Verification:} CAS (verify Hopf algebra axioms)
\end{problembox}

\subsection*{A.3: Counterexample --- Locally Finite Type vs.\ Finite Type (Level 2)}

\begin{problembox}[title={\ecagid{0043} --- Basic Concepts}]
\textbf{Problem:} State the precise difference between ``locally of finite type'' and ``of finite type'' for a morphism of schemes. Give an example of a morphism that is locally of finite type but not of finite type, and identify which condition fails.

\medskip
\textbf{Solution (abbreviated):}
A morphism $f: X \to Y$ is \emph{locally of finite type} if for every $x \in X$, there exist affine neighborhoods $U = \Spec(A)$ of $f(x)$ and $V = \Spec(B) \subseteq f^{-1}(U)$ of $x$ such that $B$ is a finitely generated $A$-algebra. It is \emph{of finite type} if additionally $f$ is quasi-compact.

The morphism $\coprod_{i \in I} \mathbf{A}^1_k \to \mathbf{A}^1_k$ (identity on each component, $I$ infinite) is locally of finite type but not of finite type, because the source is not quasi-compact.

\medskip
\textbf{Type:} Counterexample \quad \textbf{Difficulty:} 2 \quad \textbf{Verification:} CAS
\end{problembox}

\subsection*{A.4: Example --- \'Etale Cover of Nodal Cubic (Level 2)}

\begin{problembox}[title={\ecagid{0158} --- Basic Concepts}]
\textbf{Problem:} Let $k$ be an algebraically closed field with $\chr(k) \neq 2$, and let $Y = \Spec(k[x,y]/(y^2 - x^2(x+1)))$ be the nodal cubic. Construct a degree-2 \'etale morphism $f: X \to Y$ and compute the fibers over the node and a smooth point.

\medskip
\textbf{Solution (abbreviated):}
Define $X = \Spec(k[s,t]/(t^2 - (s^2-1)^2))$ with $f$ given by $x \mapsto s^2 - 1$, $y \mapsto st$.
The fiber over the node $(0,0)$ is $\Spec(k[s,t]/(s^2-1, t^2-(s^2-1)^2, st)) = \{(1,0), (-1,0)\}$, a reduced scheme of two points.
At smooth points, $f$ is \'etale of degree~2 (one checks that $\Omega_{X/Y} = 0$ away from the node).

\medskip
\textbf{Type:} Example \quad \textbf{Difficulty:} 2 \quad \textbf{Verification:} CAS (verify \'etale condition, compute fibers)
\end{problembox}

\subsection*{A.5: Computation --- Horrocks--Mumford Bundle (Level 2)}

\begin{problembox}[title={\ecagid{0230} --- Computations}]
\textbf{Problem:} Let $\mathscr{F}$ be the Horrocks--Mumford bundle on $\mathbf{P}^4$ with $c_1(\mathscr{F}) = 5h$ and $c_2(\mathscr{F}) = 10h^2$. Using HRR, compute $\chi(\mathscr{F}(n-5))$ as a closed-form polynomial in $n$.

\medskip
\textbf{Solution (abbreviated):}
The Chern character of $\mathscr{F}(m)$ (where $m = n-5$) is computed from $c_1$ and $c_2$ using the splitting principle, then paired with the Todd class of $\mathbf{P}^4$.
The result is a degree-4 polynomial in $n$, which at $n=5$ yields $\chi(\mathscr{F}) = h^0 - h^1 = 4 - 2 = 2$, consistent with known cohomology.

\medskip
\textbf{Type:} Computation \quad \textbf{Difficulty:} 2 \quad \textbf{Verification:} Exact match + CAS
\end{problembox}

\subsection*{A.6: Example --- Proper Birational to Normal (Level 2)}

\begin{problembox}[title={\ecagid{0019} --- Basic Concepts}]
\textbf{Problem:} Let $f: X \to Y$ be a proper birational morphism with $Y$ normal. Prove that $f_* \mathcal{O}_X \cong \mathcal{O}_Y$. Which property of normal schemes is used?

\medskip
\textbf{Solution (abbreviated):}
The structure map $\mathcal{O}_Y \to f_* \mathcal{O}_X$ is injective (since $f$ is birational and dominant).
For surjectivity, any section $s \in (f_* \mathcal{O}_X)(U) = \mathcal{O}_X(f^{-1}(U))$ restricts to a rational function on $U$ that is regular on the dense open $f^{-1}(U \cap V)$ where $f$ is an isomorphism.
Since $Y$ is normal, the algebraic Hartogs lemma (Serre's $R_1 + S_2$ criterion) extends $s$ to a regular function on $U$.

\medskip
\textbf{Type:} Example \quad \textbf{Difficulty:} 2 \quad \textbf{Verification:} Rubric
\end{problembox}

\section{Evaluation Prompt and Scoring Rubric}
\label{app:prompt}

\subsection*{B.1: Full Evaluation Prompt}

\begin{lstlisting}
You are an expert algebraic geometer with deep
knowledge of scheme theory, sheaf cohomology,
moduli spaces, and computational algebraic geometry.

PROBLEM:
{problem}

INSTRUCTIONS:
Provide a CONCRETE construction with full
mathematical justification. Your response should
include:

1. **The explicit mathematical object**: State the
   object clearly (e.g., a specific polynomial,
   ring, scheme, morphism, sheaf, or cohomology
   class).

2. **Verification of all required properties**: For
   each property mentioned in the problem, provide
   a rigorous argument that your object satisfies
   (or, for counterexamples, violates) it.

3. **Key insight or theorem used**: Identify the main
   mathematical idea underlying your construction.

FORMAT:
- State your answer at the very beginning, clearly
  labeled.
- Then provide the detailed justification.
- Use standard LaTeX notation for mathematical
  expressions.
- If the problem asks for a counterexample,
  explicitly state which hypothesis fails and why.
\end{lstlisting}

\subsection*{B.2: Scoring Rubric}

\begin{description}[nosep]
  \item[Score 1.0 (Correct):] \hfill
  \begin{itemize}[nosep]
    \item (C1) The constructed object is well-defined and belongs to the correct mathematical category.
    \item (C2) All properties stated in the problem are rigorously verified (or, for counterexamples, the target property is rigorously shown to fail).
    \item (C3) The answer is non-trivial and non-degenerate.
    \item (C4) The mathematical reasoning is sound and references to standard results are accurate.
  \end{itemize}
  \item[Score 0.5 (Partial):] \hfill
  \begin{itemize}[nosep]
    \item The object is valid (C1 satisfied) but one or more of C2--C4 is partially satisfied.
    \item Typical cases: correct object with incomplete verification, or correct approach with a computational error that does not invalidate the construction.
  \end{itemize}
  \item[Score 0.0 (Incorrect):] \hfill
  \begin{itemize}[nosep]
    \item The object is invalid (C1 fails), or the properties are not satisfied (C2 fails fundamentally).
    \item Includes fabricated objects (F1), reversed properties (F2), and reliance on non-existent theorems (F5).
  \end{itemize}
\end{description}

\section{Dataset Schema and Access}
\label{app:schema}

\subsection*{C.1: JSON Schema}

Each entry in the \ecagbench{} dataset (\texttt{benchmark/data/ecag\_bench.jsonl}) conforms to the following schema:

\begin{lstlisting}
{
  "id": "ecag-XXXX",
  "problem": "Problem statement in LaTeX",
  "solution": "Full solution text",
  "answer": "Closed-form answer or null",
  "answer_type": "example|counterexample|
                  closed_form|proof|true_false",
  "difficulty": 1|2|3,
  "chapter": "basic_concepts|cohomology|
              computations|moduli|relative",
  "topic": "...",
  "subtopic": "...",
  "question_type": "example|counterexample|
                    computation|proof|true_false",
  "tags": ["..."],
  "source": "Original|Hartshorne X.Y.Z|...",
  "prerequisites": ["..."],
  "msc_codes": ["14XXX"],
  "verification_strategy": "exact|cas|rubric|hybrid",
  "verification_script": "path or null",
  "accepts_multiple": true|false,
  "is_stub": false,
  "latex_source": "path to .tex file"
}
\end{lstlisting}

Difficulty levels: 1 = routine (follows from definitions), 2 = non-trivial (requires significant theorem), 3 = deep (combines multiple areas or requires specialized knowledge).

\subsection*{C.2: Access}

\begin{itemize}[nosep]
  \item \textbf{GitHub:} \url{https://github.com/Waerden001/Examples_and_Counterexamples_in_Elementary_Algebraic_Geometry}
  \item \textbf{Website:} The accompanying textbook is available as a searchable website via GitHub Pages.
  \item \textbf{License:} CC-BY-4.0
\end{itemize}

\section{Case Study Full Outputs}
\label{app:case-studies}

This appendix provides representative model outputs from the pilot evaluation discussed in Section~\ref{sec:experiments}.
All outputs are from Claude Sonnet~4.6 via \texttt{claude -p}.

\subsection*{D.1: Successful Counterexample (\ecagid{0043})}

\textbf{Problem:} State the precise difference between ``locally of finite type'' and ``of finite type'' for a morphism of schemes. Give an example of a morphism that is locally of finite type but not of finite type, and identify which condition fails.

\textbf{Type:} Counterexample \quad \textbf{Difficulty:} Level~2 \quad \textbf{Score:} 1.0

\textbf{Model output (abbreviated):}

\begin{quote}
\small
\textbf{Answer:} The precise difference lies in \textbf{quasi-compactness}: finite type = locally of finite type + quasi-compact.

\textbf{Counterexample:} Let $X = \coprod_{n=0}^{\infty} \mathbb{A}^n_k$, $Y = \Spec(k)$, $f: X \to Y$ the structure morphism. Then $f$ is locally of finite type but not of finite type; the failing condition is quasi-compactness.

\textbf{Verification --- locally of finite type:} Take the affine open cover $\{V = \Spec(k)\}$ of $Y$. Then $f^{-1}(V) = X = \coprod_{n \geq 0} \Spec(k[t_1,\ldots,t_n])$. Each $k[t_1,\ldots,t_n]$ is a finitely generated $k$-algebra, so $f$ is locally of finite type.

\textbf{Verification --- not quasi-compact:} The open cover $\{U_n\}_{n \geq 0}$ where $U_n = \Spec(k[t_1,\ldots,t_n])$ has no finite subcover, since each $U_n$ is a distinct connected component.

\textbf{Key insight:} The distinction between ``locally'' and ``globally'' properties of morphisms typically involves a quasi-compactness condition.
\end{quote}

The model correctly identified quasi-compactness as the distinguishing property, constructed a standard separating example, and verified both the positive and negative conditions rigorously.

\subsection*{D.2: Property Confusion (\ecagid{0343}, F2)}

\textbf{Problem:} In the setup above, verify that $g \cdot P = Q$, $g \cdot C = D$, and that the fixed locus of the $G$-action on $\mathbb{P}^{3}$ is exactly $l_{1} \cup l_{2}$. Compute the invariant ring $k[x,y,z,w]^{G}$ and verify the relation $v_{0}v_{1} = v_{2}^{2}$.

\textbf{Type:} Counterexample \quad \textbf{Difficulty:} Level~3 \quad \textbf{Score:} 0.5 \quad \textbf{Failure mode:} F2

\textbf{Model output (abbreviated):}

\begin{quote}
\small
\textbf{Setup:} Let $G = \mathbb{Z}/2\mathbb{Z} = \{e, g\}$, with $g$ acting on $\mathbb{P}^3_k$ by:
$g \cdot [x:y:z:w] := [-x:-y:z:w]$
corresponding to $\sigma = \operatorname{diag}(-1,-1,1,1)$.

\textbf{Verification $g \cdot P = Q$:}
$g \cdot P = g \cdot [1:0:-1:0] = [-1:0:-1:0] = [1:0:1:0] = Q$. $\checkmark$

\textbf{Invariant ring:} The degree-$d$ invariants of $k[x,y,z,w]$ under this action consist of monomials $x^a y^b z^c w^d$ with $a+b$ even.  Setting $v_0 = x^2, v_1 = y^2, v_2 = xy, v_3 = z, v_4 = w$, the invariant ring is $k[v_0, v_1, v_2, v_3, v_4]/(v_0 v_1 - v_2^2)$. $\checkmark$
\end{quote}

\textbf{Analysis:} The model substituted a different group action ($g \cdot [x:y:z:w] = [-x:-y:z:w]$) from the one specified in the problem setup (which is $g \cdot [x:y:z:w] = [y:x:w:z]$, a transposition). While the algebraic computations for the substituted action were internally consistent, the mismatch with the intended construction means the results do not answer the stated question. This exemplifies F2 (property confusion): the model demonstrates algebraic competence but applies it to the wrong mathematical object.

\subsection*{D.3: Timeout on Deep Problem (\ecagid{0363})}

\textbf{Problem:} Using genus-0 Gromov--Witten invariants and Schubert calculus on $\operatorname{Gr}(2, n+1)$, compute $\langle H^i | H^j \rangle_{0,d}$ and verify the quantum product formula.

\textbf{Type:} Example \quad \textbf{Difficulty:} Level~3 \quad \textbf{Score:} 0.0 (timeout)

The model timed out after 600 seconds without producing any response. This problem requires synthesizing Gromov--Witten theory, Schubert calculus, and quantum cohomology---a combination of techniques at the frontier of the represented material. The timeout pattern was common at Level~3, where 45.5\% of problems caused the model to fail silently, suggesting that the model's inability to produce even a partial response on such problems reflects a fundamental gap in its ability to combine deep techniques from multiple areas of algebraic geometry.
